% Options for packages loaded elsewhere
\PassOptionsToPackage{unicode}{hyperref}
\PassOptionsToPackage{hyphens}{url}
\PassOptionsToPackage{dvipsnames,svgnames,x11names}{xcolor}
\documentclass[
  10pt,
  fontset=fandol]{ctexart}
\usepackage{xcolor}
\usepackage[margin=16mm]{geometry}
\usepackage{amsmath,amssymb}
\setcounter{secnumdepth}{-\maxdimen} % remove section numbering
\usepackage{iftex}
\ifPDFTeX
  \usepackage[T1]{fontenc}
  \usepackage[utf8]{inputenc}
  \usepackage{textcomp} % provide euro and other symbols
\else % if luatex or xetex
  \usepackage{unicode-math} % this also loads fontspec
  \defaultfontfeatures{Scale=MatchLowercase}
  \defaultfontfeatures[\rmfamily]{Ligatures=TeX,Scale=1}
\fi
\usepackage{lmodern}
\ifPDFTeX\else
  % xetex/luatex font selection
\fi
% Use upquote if available, for straight quotes in verbatim environments
\IfFileExists{upquote.sty}{\usepackage{upquote}}{}
\IfFileExists{microtype.sty}{% use microtype if available
  \usepackage[]{microtype}
  \UseMicrotypeSet[protrusion]{basicmath} % disable protrusion for tt fonts
}{}
\makeatletter
\@ifundefined{KOMAClassName}{% if non-KOMA class
  \IfFileExists{parskip.sty}{%
    \usepackage{parskip}
  }{% else
    \setlength{\parindent}{0pt}
    \setlength{\parskip}{6pt plus 2pt minus 1pt}}
}{% if KOMA class
  \KOMAoptions{parskip=half}}
\makeatother
\setlength{\emergencystretch}{3em} % prevent overfull lines
\providecommand{\tightlist}{%
  \setlength{\itemsep}{0pt}\setlength{\parskip}{0pt}}
\usepackage{amsmath,amssymb,mathtools,needspace}
\xeCJKDeclareCharClass{CJK}{"2032,"0400->"04FF,"2160->"216B,"2460->"2473,"25A0,"25C7,"25CB,"25CF}
\IfFontExistsTF{Arial Unicode MS}{\xeCJKsetup{AutoFallBack=true}\setCJKfallbackfamilyfont{\CJKrmdefault}{Arial Unicode MS}}{}
\setcounter{secnumdepth}{0}
\setlength{\emergencystretch}{2em}
\allowdisplaybreaks
\usepackage{bookmark}
\IfFileExists{xurl.sty}{\usepackage{xurl}}{} % add URL line breaks if available
\urlstyle{same}
\hypersetup{
  pdftitle={Hamming 格式},
  colorlinks=true,
  linkcolor={Maroon},
  filecolor={Maroon},
  citecolor={Blue},
  urlcolor={Blue},
  pdfcreator={LaTeX via pandoc}}

\title{Hamming 格式}
\author{}
\date{}

\begin{document}
\maketitle

\subsubsection{5.5.7 Hamming 格式}\label{hamming-ux683cux5f0f}

上述预测-校正系统的缺点是稳定性差。为了改善稳定性，重新挑选校正公式。为此考察式（5.5.21）中不显含 \(y'_{n-2}\) 的一类格式

\[
y_{n+1}=\alpha_0y_n+\alpha_1y_{n-1}+\alpha_2y_{n-2}+h(\beta_{-1}y'_{n+1}+\beta_0y'_n+\beta_1y'_{n-1}),
\tag{5.5.28}
\]

按条件（5.5.22），形如式（5.5.28）的四阶格式含有一个自由度------譬如可取 \(\alpha_1\) 为独立参数。若取 \(\alpha_1=1\) 即得 Simpson 格式（5.5.27）。

Hamming 用 \(\alpha_1\) 的不同数值进行试验，发现当 \(\alpha_1=0\) 时格式的稳定性能较好，即

\[
y_{n+1}=\frac18(9y_n-y_{n-2})+\frac{3h}{8}(y'_{n+1}+2y'_n-y'_{n-1}),
\tag{5.5.29}
\]

其局部截断误差为

\[
y(x_{n+1})-y_{n+1}\approx-\frac{h^5}{40}y_n^{(5)}.
\tag{5.5.30}
\]

Hamming 格式（5.5.29）不能用数值积分方法推导出来。

用 Milne 格式（5.5.24）与 Hamming 格式（5.5.29）相匹配，并利用误差公式（5.5.25）、（5.5.30）改进计算结果，即可建立下列预测-校正系统：

预测

\[
p_{n+1}=y_{n-3}+\frac{4h}{3}(2y'_n-y'_{n-1}+2y'_{n-2}),
\]

改进

\[
m_{n+1}=p_{n+1}-\frac{112}{121}(p_n-c_n),
\]

计算

\[
m'_{n+1}=f(x_{n+1},m_{n+1}),
\]

校正

\[
c_{n+1}=\frac18(9y_n-y_{n-2})+\frac{3h}{8}(m'_{n+1}+2y'_n-y'_{n-1}),
\]

改进

\[
y_{n+1}=c_{n+1}+\frac9{121}(p_{n+1}-c_{n+1}),
\]

\href{../../source-images/pdf-145.jpeg}{原页}

计算

\[
y'_{n+1}=f(x_{n+1},y_{n+1}).
\]

\end{document}
